\documentclass[12pt]{article}
\usepackage{amsmath}
\usepackage{amsfonts}
\usepackage{float}
\usepackage{fancyhdr}
\usepackage{graphicx}
\usepackage{setspace}
\usepackage[colorlinks=true,linkcolor=blue, citecolor=red]{hyperref}
\usepackage{url}
\usepackage[top=.75in, left=.75in, right=.75in, bottom=1in]{geometry}

% For algorithm
\usepackage{algorithm}
\usepackage{algpseudocode}

% ============ CODE ============
\usepackage{listings}
\usepackage{xcolor}
\definecolor{codegreen}{rgb}{0,0.6,0}
\definecolor{codegray}{rgb}{0.5,0.5,0.5}
\definecolor{codepurple}{rgb}{0.58,0,0.82}
\definecolor{backcolour}{rgb}{0.95,0.95,0.92}

% Styling for the code.
\lstdefinestyle{mystyle}{
    backgroundcolor=\color{backcolour},   
    commentstyle=\color{codegreen},
    keywordstyle=\color{magenta},
    numberstyle=\tiny\color{codegray},
    stringstyle=\color{codepurple},
    basicstyle=\ttfamily\footnotesize,
    breakatwhitespace=false,         
    breaklines=true,                 
    captionpos=b,                    
    keepspaces=true,                 
    numbers=left,                    
    numbersep=5pt,                  
    showspaces=false,                
    showstringspaces=false,
    showtabs=false,                  
    tabsize=2
}
\lstset{style=mystyle}

% Table packages and Unicode / Font support
\usepackage{array}
\usepackage{longtable}
\usepackage{enumitem}
\usepackage{booktabs}
\usepackage[utf8]{inputenc}
\DeclareUnicodeCharacter{2264}{\ensuremath{\le}}
\DeclareUnicodeCharacter{2194}{\ensuremath{\leftrightarrow}}
\usepackage[T1]{fontenc}
\usepackage{lmodern}
\graphicspath{{../../../../templates/}{../../prototype/}{figures/}{./}}
\setlength{\emergencystretch}{3em}
\setlist[enumerate]{label=\arabic*.,leftmargin=*}
\setlist[itemize]{leftmargin=*}

% ============ HEADER AND FOOTER ============
\setlength{\headheight}{29.43912pt}

\pagestyle{fancy}
\fancyfoot{}
\fancyfoot[R]{Page \thepage}

\setlength{\parindent}{0pt}
\renewcommand{\baselinestretch}{1.5}

\newcommand{\coursename}{Introduction of Software Engineering}
\newcommand{\reportname}{Studitopia}
\newcommand{\reporttitle}{\href{https://github.com/RoundFloor/Studitopia.git}{Studitopia}}

\newcommand{\studentname}{%
Nguyen Thanh An (24125024)\\
Nguyen Chi Bao (24125025)\\
Tran Nguyen Khanh Binh (24125026)\\
Tran Canh Anh Tuan (24125048)\\
Nguyen Le Quoc Huy (24125100)
}
\newcommand{\teachername}{Ho Tuan Thanh, Nguyen Thi Minh Tuyen}

% Header
\lhead{\reporttitle}
\rhead{
University of Science - VNUHCM\\
\coursename
}

% Footer
\newcommand{\leftfooter}{}
\lfoot{\leftfooter}

% ============ DOCUMENT ============
\begin{document}

\pagenumbering{roman}
\begin{titlepage}
\newcommand{\HRule}{\rule{\linewidth}{0.5mm}}
\centering

\textsc{\LARGE VNU-HCM}\\[0.5cm]
\textsc{\Large University of Science}\\[0.5cm]
\textsc{\large Faculty of Information Technology}\\[0.5cm]
\textsc{Computer Science}\\[0.5cm]

\HRule \\[0.4cm]
{ 
\huge{\bfseries{\reporttitle}}\\[0.5cm]
\large{\bfseries{UC-23 --- Receive Resource Recommendations}}
}\\[0.4cm]
\HRule \\[0.5cm]

\textbf{\large Subject: \coursename}\\[0.5cm]

\begin{minipage}[t]{0.55\textwidth}
\begin{flushleft} \large
\emph{Students:}\\
\studentname
\end{flushleft}
\end{minipage}
~
\begin{minipage}[t]{0.4\textwidth}
\begin{flushright} \large
\emph{Supervising Teachers:} \\
\teachername
\end{flushright}
\end{minipage}\\[1cm]

{\large \today}\\[1cm]

\includegraphics[scale=.20]{img/hcmus-logo.png}\\[1cm] 

\vfill
\end{titlepage}	

\tableofcontents
\pagebreak

\pagenumbering{arabic}
\setcounter{page}{1}

\hypertarget{sec:uc-23-receive-resource-recommendations}{}
\addcontentsline{toc}{section}{UC-23 --- Receive Resource Recommendations}
\section*{UC-23 --- Receive Resource Recommendations}
\par\medskip

\textbf{Version:} \texttt{1.0}
\par\medskip

\textbf{Performed by:} [Nguyen Le Quoc Huy] | \textbf{Edited by:} [Nguyen Le Quoc Huy]
\par\medskip

\begingroup
\small
\renewcommand{\baselinestretch}{1.0}\selectfont
\renewcommand{\arraystretch}{1.15}
\setlength{\tabcolsep}{6pt}
\begin{longtable}{|>{\raggedright\arraybackslash}p{0.09\textwidth}|>{\raggedright\arraybackslash}p{0.84\textwidth}|}
\hline
\textbf{Field} & \textbf{Value} \\ \hline
\endhead
Use Case ID & UC-23 \\ \hline
Name & Receive Resource Recommendations \\ \hline
Primary Actor & University Student \\ \hline
Supporting Actors & None \\ \hline
Feature Group & Learning Resource Library \\ \hline
Priority & Not explicitly stated \\ \hline
Related User Stories & US-26 \\ \hline
Related Requirements & FR-26.1 to FR-26.4; AC-26.1 to AC-26.3; FR-09.4; FR-09.5; FR-25.5; FR-27.4; NFR-PERF-1; NFR-SCAL-4; NFR-SEC-1; NFR-SEC-6; NFR-PRIV-5; NFR-PRIV-7; NFR-USE-4 \\ \hline
Source Documents & \texttt{docs/requirements/proposals/userstory.md}; \texttt{docs/requirements/proposals/functional\_requirement.md}; \texttt{docs/requirements/proposals/library.md}; \texttt{docs/requirements/proposals/create\_post.md}; \texttt{docs/requirements/proposals/API.md}; \texttt{docs/requirements/proposals/response\_code.md}; \texttt{docs/requirements/proposals/non\_functional\_requirement.md}; \texttt{docs/requirements/proposals/feature\_flow.md}; \texttt{docs/requirements/proposals/screen\_flow.md} \\ \hline
\end{longtable}
\endgroup
\par\medskip

\hypertarget{sec:part-a-use-case-model}{}
\addcontentsline{toc}{subsection}{Part A --- Use Case Model}
\subsection*{Part A --- Use Case Model}
\par\medskip

\hypertarget{sec:a1-focused-use-case-diagram}{}
\addcontentsline{toc}{subsubsection}{A1. Focused Use Case Diagram}
\subsubsection*{A1. Focused Use Case Diagram}
\par\medskip

\begin{center}
\includegraphics[width=0.85\textwidth,height=0.75\textheight,keepaspectratio]{"figures/UC-23-resource-recommendations_1.png"}
\end{center}
\par\medskip

\textbf{Diagram notation.} Person shape = human actor; cloud = supporting system; blue = selected use case; solid line = association; labelled solid arrow = generalization; dashed arrow = include or extend.
\par\medskip

\hypertarget{sec:a2-actor-register}{}
\addcontentsline{toc}{subsubsection}{A2. Actor Register}
\subsubsection*{A2. Actor Register}
\par\medskip

\begingroup
\small
\renewcommand{\baselinestretch}{1.0}\selectfont
\renewcommand{\arraystretch}{1.15}
\setlength{\tabcolsep}{6pt}
\begin{longtable}{|>{\raggedright\arraybackslash}p{0.11\textwidth}|>{\raggedright\arraybackslash}p{0.09\textwidth}|>{\raggedright\arraybackslash}p{0.71\textwidth}|}
\hline
\textbf{Actor} & \textbf{Type} & \textbf{Description} \\ \hline
\endhead
University Student & Primary & Authenticated student who receives resource suggestions based on subjects, courses, interests, and academic context. \\ \hline
\end{longtable}
\endgroup
\par\medskip

\hypertarget{sec:a3-relationships}{}
\addcontentsline{toc}{subsubsection}{A3. Relationships}
\subsubsection*{A3. Relationships}
\par\medskip

\hypertarget{sec:a3-1-associations}{}
\paragraph*{A3.1 Associations}
\par\medskip

\begingroup
\small
\renewcommand{\baselinestretch}{1.0}\selectfont
\renewcommand{\arraystretch}{1.15}
\setlength{\tabcolsep}{6pt}
\begin{longtable}{|>{\raggedright\arraybackslash}p{0.13\textwidth}|>{\raggedright\arraybackslash}p{0.27\textwidth}|>{\raggedright\arraybackslash}p{0.09\textwidth}|>{\raggedright\arraybackslash}p{0.41\textwidth}|}
\hline
\textbf{Actor} & \textbf{Use Case} & \textbf{Type} & \textbf{Description} \\ \hline
\endhead
University Student & UC-23 Receive Resource Recommendations & Primary & Views and saves recommended books and learning resources. \\ \hline
\end{longtable}
\endgroup
\par\medskip

\hypertarget{sec:a3-2-generalization-and-specialization}{}
\paragraph*{A3.2 Generalization and Specialization}
\par\medskip

\begingroup
\small
\renewcommand{\baselinestretch}{1.0}\selectfont
\renewcommand{\arraystretch}{1.15}
\setlength{\tabcolsep}{6pt}
\begin{longtable}{|>{\raggedright\arraybackslash}p{0.09\textwidth}|>{\raggedright\arraybackslash}p{0.09\textwidth}|>{\raggedright\arraybackslash}p{0.19\textwidth}|>{\raggedright\arraybackslash}p{0.54\textwidth}|}
\hline
\textbf{Child} & \textbf{Parent} & \textbf{Type} & \textbf{Inherited Behavior or Association} \\ \hline
\endhead
None & None & No generalization or specialization & The proposals describe new users as a recommendation condition, not as a distinct actor for US-26. \\ \hline
\end{longtable}
\endgroup
\par\medskip

\hypertarget{sec:a3-3-include}{}
\paragraph*{A3.3 Include}
\par\medskip

\begingroup
\small
\renewcommand{\baselinestretch}{1.0}\selectfont
\renewcommand{\arraystretch}{1.15}
\setlength{\tabcolsep}{6pt}
\begin{longtable}{|>{\raggedright\arraybackslash}p{0.23\textwidth}|>{\raggedright\arraybackslash}p{0.30\textwidth}|>{\raggedright\arraybackslash}p{0.40\textwidth}|}
\hline
\textbf{Base Use Case} & \textbf{Included Use Case} & \textbf{Meaning} \\ \hline
\endhead
None & None & No include relationship & The proposals do not name an included use case for US-26. \\ \hline
\end{longtable}
\endgroup
\par\medskip

\hypertarget{sec:a3-4-extend}{}
\paragraph*{A3.4 Extend}
\par\medskip

\begingroup
\small
\renewcommand{\baselinestretch}{1.0}\selectfont
\renewcommand{\arraystretch}{1.15}
\setlength{\tabcolsep}{6pt}
\begin{longtable}{|>{\raggedright\arraybackslash}p{0.15\textwidth}|>{\raggedright\arraybackslash}p{0.11\textwidth}|>{\raggedright\arraybackslash}p{0.18\textwidth}|>{\raggedright\arraybackslash}p{0.47\textwidth}|}
\hline
\textbf{Extending Use Case} & \textbf{Base Use Case} & \textbf{Condition} & \textbf{Added Behavior} \\ \hline
\endhead
None & None & No extend relationship & The proposals do not name an extending use case for US-26. \\ \hline
\end{longtable}
\endgroup
\par\medskip

\hypertarget{sec:a4-model-notes}{}
\addcontentsline{toc}{subsubsection}{A4. Model Notes}
\subsubsection*{A4. Model Notes}
\par\medskip

\begin{enumerate}
\item The scope is personalized and cold-start resource suggestions, their display, and saving recommended resources to a personal list.
\item Searching by query, publishing resources, and deleting published resources are outside this use case.
\item Dismissal and refresh controls are library.md UI recommendations; they are not modeled as new extending use cases.
\item No distinct included, extending, or specialized use case is named for US-26, so no I-xx, E-xx, or G-xx relationship ID is introduced.
\item The exact API endpoint, route, recommendation count, and priority are not specified in proposals.
\end{enumerate}
\par\medskip

\hypertarget{sec:part-b-use-case-specification}{}
\addcontentsline{toc}{subsection}{Part B --- Use Case Specification}
\subsection*{Part B --- Use Case Specification}
\par\medskip

\hypertarget{sec:b1-brief-description}{}
\addcontentsline{toc}{subsubsection}{B1. Brief Description}
\subsubsection*{B1. Brief Description}
\par\medskip

A university student receives recommendations for books and learning resources so that the student can discover useful materials related to subjects. The system recommends resources related to the student's subjects, courses, and interests, displays them on the home or library page, and allows the student to save recommended resources to a personal list.
\par\medskip

For new users without historical activity, the system provides curated foundational resources. The library project frames discovery around \textbf{For you}, \textbf{Top reading}, \textbf{Course progress}, and \textbf{Starter set}, with user-relevance and course-progress recommendations drawn from the feature flows.
\par\medskip

\hypertarget{sec:b2-preconditions}{}
\addcontentsline{toc}{subsubsection}{B2. Preconditions}
\subsubsection*{B2. Preconditions}
\par\medskip

\begin{enumerate}
\item The student is authenticated before accessing protected user resources (\texttt{NFR-SEC-1}).
\item The Learning Resource Library is reachable through \texttt{Home -$>$ Library Group} in \texttt{screen\_flow.md}.
\item The student's preferences, courses, subjects, and interests are available where the student has provided them (\texttt{FR-09.1}, \texttt{FR-09.2}, \texttt{FR-09.5}; \texttt{AC-09.1}, \texttt{AC-09.2}, \texttt{AC-09.5}).
\item Published resources are available for recommendation after publication (\texttt{FR-27.4}, \texttt{AC-27.1}).
\item For a new user, limited historical activity is treated as the cold-start condition (\texttt{FR-26.4}, \texttt{AC-26.3}).
\end{enumerate}
\par\medskip

\hypertarget{sec:b3-postconditions}{}
\addcontentsline{toc}{subsubsection}{B3. Postconditions}
\subsubsection*{B3. Postconditions}
\par\medskip

\begin{enumerate}
\item \textbf{Success:} Recommended resources are displayed on the home or library page.
\item \textbf{Success:} Suggestions reflect the student's preferences and courses, or curated foundational resources for a new user.
\item \textbf{Success:} A recommended resource selected for saving is added to the student's personal list.
\item \textbf{Failure:} No recommendation is presented as personalized when the system has no basis, and the specified cold-start or load-error feedback is shown.
\end{enumerate}
\par\medskip

\hypertarget{sec:b4-main-success-scenario-basic-flow}{}
\addcontentsline{toc}{subsubsection}{B4. Main Success Scenario (Basic Flow)}
\subsubsection*{B4. Main Success Scenario (Basic Flow)}
\par\medskip

\begin{enumerate}
\item The student opens the library recommendation area or the home-page recommendation area.
\item The system displays \texttt{RecommendationPanel} with heading \textbf{Recommended for your studies}.
\item The student selects a recommendation segment such as \textbf{For you}, \textbf{Top reading}, \textbf{Course progress}, or \textbf{Starter set}.
\item The system recommends books and learning resources related to the student's subjects, courses, and interests (\texttt{FR-26.1}, \texttt{AC-26.1}).
\item The system displays each recommendation with \texttt{RecommendationReasonText}, for example “Recommended because you are enrolled in CS101.”
\item The student reviews the recommendation card with title, type, tags, rating, and access information.
\item The student selects \textbf{Save} through \texttt{SaveRecommendationButton}.
\item The system saves the recommended resource to the student's personal list (\texttt{FR-26.3}, \texttt{AC-26.2}).
\item The saved item becomes available through \texttt{SavedResourceList}.
\end{enumerate}
\par\medskip

\hypertarget{sec:b5-alternative-flows}{}
\addcontentsline{toc}{subsubsection}{B5. Alternative Flows}
\subsubsection*{B5. Alternative Flows}
\par\medskip

\hypertarget{sec:a1-home-page-display}{}
\paragraph*{A1. Home-page display}
\par\medskip

The recommendation set is shown on the student's home page instead of the library page (\texttt{FR-26.2}). The same save action adds the resource to the personal list.
\par\medskip

\hypertarget{sec:a2-top-reading-or-course-progress}{}
\paragraph*{A2. Top reading or course progress}
\par\medskip

The student selects \textbf{Top reading} or \textbf{Course progress}. The system displays popular or course-connected resources with rating and access information; these are library.md recommendation segments derived from \texttt{feature\_flow.md} and the course-progress flow.
\par\medskip

\hypertarget{sec:a3-cold-start-starter-resources}{}
\paragraph*{A3. Cold-start starter resources}
\par\medskip

The student is new or has little historical activity. The system displays \texttt{CuratedStarterResourceList} with foundational resources based on selected subjects (\texttt{FR-26.4}, \texttt{AC-26.3}).
\par\medskip

\hypertarget{sec:a4-dismiss-or-refresh-suggestions}{}
\paragraph*{A4. Dismiss or refresh suggestions}
\par\medskip

The student selects \textbf{Not relevant} to hide an item or \textbf{Refresh suggestions} to request a new set. These controls are explicit library.md recommendations and are not confirmed FRs for US-26.
\par\medskip

\hypertarget{sec:b6-exception-flows}{}
\addcontentsline{toc}{subsubsection}{B6. Exception Flows}
\subsubsection*{B6. Exception Flows}
\par\medskip

\hypertarget{sec:e1-recommendation-load-failure}{}
\paragraph*{E1. Recommendation load failure}
\par\medskip

The recommendation set cannot be loaded. The system shows “Unable to load resources.” with \textbf{Retry} through \texttt{LibraryErrorBanner}.
\par\medskip

\hypertarget{sec:e2-recommendation-reason-unavailable}{}
\paragraph*{E2. Recommendation reason unavailable}
\par\medskip

No explanation basis is available for a suggested resource. The system omits \texttt{RecommendationReasonText} rather than displaying an invented reason.
\par\medskip

\hypertarget{sec:e3-save-failure}{}
\paragraph*{E3. Save failure}
\par\medskip

The save action does not complete. The system keeps the recommendation available for another attempt; the exact save-failure message is Not specified in proposals.
\par\medskip

\hypertarget{sec:b7-business-rules}{}
\addcontentsline{toc}{subsubsection}{B7. Business Rules}
\subsubsection*{B7. Business Rules}
\par\medskip

\begin{enumerate}
\item The system recommends books or learning resources related to the user's subjects, courses, and interests (\texttt{FR-26.1}, \texttt{AC-26.1}).
\item The system displays recommended resources on the user's home or library page (\texttt{FR-26.2}).
\item Users can save recommended resources to a personal list (\texttt{FR-26.3}, \texttt{AC-26.2}).
\item New users without historical activity receive curated foundational resources (\texttt{FR-26.4}, \texttt{AC-26.3}).
\item Learning preferences are used to recommend materials (\texttt{FR-09.4}), and preferences can be updated at any time (\texttt{FR-09.5}, \texttt{AC-09.5}).
\item The library displays resources related to the user's university, major, and courses (\texttt{FR-25.5}, \texttt{AC-25.4}).
\item Access to recommended resources follows owner privacy settings (\texttt{FR-27.5}, \texttt{NFR-PRIV-5}, \texttt{NFR-PRIV-7}).
\end{enumerate}
\par\medskip

\hypertarget{sec:b8-special-requirements}{}
\addcontentsline{toc}{subsubsection}{B8. Special Requirements}
\subsubsection*{B8. Special Requirements}
\par\medskip

\begingroup
\small
\renewcommand{\baselinestretch}{1.0}\selectfont
\renewcommand{\arraystretch}{1.15}
\setlength{\tabcolsep}{6pt}
\begin{longtable}{|>{\raggedright\arraybackslash}p{0.10\textwidth}|>{\raggedright\arraybackslash}p{0.83\textwidth}|}
\hline
\textbf{ID} & \textbf{Requirement} \\ \hline
\endhead
NFR-PERF-1 & The system shall provide responsive user interactions with a response time of ≤ 2 seconds for standard user actions. \\ \hline
NFR-SCAL-4 & The system architecture shall support at least 100,000 uploaded documents and learning resources. \\ \hline
NFR-SEC-1 & The system shall require authentication before accessing protected user resources. \\ \hline
NFR-SEC-6 & The system shall prevent unauthorized access through server-side authorization checks. \\ \hline
NFR-PRIV-5 & Users shall be able to control access to documents and learning resources they share. \\ \hline
NFR-PRIV-7 & Privacy settings shall apply consistently to forum posts, comments, member lists, and search results. \\ \hline
NFR-USE-4 & Primary interface actions shall use clear labels and recognizable controls. \\ \hline
Design constraint & Recommendation segments, reason text, and dismissal/refresh controls follow \texttt{library.md} as UI recommendations; FR/AC/NFR entries are confirmed requirements. \\ \hline
\end{longtable}
\endgroup
\par\medskip

\hypertarget{sec:b9-extension-points}{}
\addcontentsline{toc}{subsubsection}{B9. Extension Points}
\subsubsection*{B9. Extension Points}
\par\medskip

\begingroup
\small
\renewcommand{\baselinestretch}{1.0}\selectfont
\renewcommand{\arraystretch}{1.15}
\setlength{\tabcolsep}{6pt}
\begin{longtable}{|>{\raggedright\arraybackslash}p{0.37\textwidth}|>{\raggedright\arraybackslash}p{0.56\textwidth}|}
\hline
\textbf{Extension Point} & \textbf{Location in Flow of Events} \\ \hline
\endhead
Home or library placement & Basic flow step 1; alternative A1 \\ \hline
Segment selection & Basic flow step 3; alternative A2 \\ \hline
Cold-start handling & Before basic flow step 4; alternative A3 \\ \hline
Dismiss or refresh & After basic flow step 5; alternative A4 \\ \hline
Recommendation load failure & Basic flow steps 2-4; exception E1 \\ \hline
Save failure & Basic flow steps 7-8; exception E3 \\ \hline
\end{longtable}
\endgroup
\par\medskip

\hypertarget{sec:b10-interface-and-data-notes}{}
\addcontentsline{toc}{subsubsection}{B10. Interface and Data Notes}
\subsubsection*{B10. Interface and Data Notes}
\par\medskip

\begin{itemize}
\item \textbf{Screen and flow path (confirmed):} \texttt{screen\_flow.md} defines \texttt{Home -$>$ Library Group}. \texttt{feature\_flow.md} defines \texttt{Library -$>$ Recommend -$>$ Top reading / User relevance}; \texttt{userstory.md} section 8 also shows \texttt{Recommend} with Top reading and User relevance, and \texttt{library.md} includes course progress.
\item \textbf{Route (not specified):} Route is Not specified in proposals; \texttt{library.md} presents \texttt{/library/recommended} and \texttt{/library/saved} as suggested states only.
\item \textbf{Recommended UI:} \texttt{RecommendationPanel} heading \textbf{Recommended for your studies}, \texttt{RecommendationSegmentControl} options \textbf{For you}, \textbf{Top reading}, \textbf{Course progress}, \textbf{Starter set}, \texttt{UserRelevanceResourceList}, \texttt{TopReadingResourceList}, \texttt{CourseProgressResourceList}, \texttt{CuratedStarterResourceList}, and \texttt{SavedResourceList}.
\item \textbf{Confirmed actions and data:} Display recommendations (\texttt{FR-26.2}, \texttt{AC-26.1}) and save them to a personal list (\texttt{FR-26.3}, \texttt{AC-26.2}). Recommended fields include \texttt{recommendationReason}, \texttt{title}, \texttt{resourceType}, \texttt{subjectTags}, \texttt{courseTags}, \texttt{rating}, \texttt{reviewCount}, \texttt{accesses}, \texttt{status}, and \texttt{publishedAt}.
\item \textbf{Exact messages and labels:} \textbf{Recommended for your studies}; \textbf{For you}; \textbf{Top reading}; \textbf{Course progress}; \textbf{Starter set}; “Recommended because you are enrolled in CS101.”; \textbf{Save}; \textbf{Not relevant}; \textbf{Refresh suggestions}; “Unable to load resources.” with \textbf{Retry}.
\item \textbf{API and codes:} Library recommendation endpoint, statuses, and error codes are Not specified in proposals. \texttt{API.md} has no standalone library recommendation endpoint, and \texttt{response\_code.md} defines no library-specific error code.
\item \textbf{Recommendation count and limits:} The number of recommendations, refresh interval, and numeric cold-start threshold are Not specified in proposals.
\item \textbf{Notification and privacy rules:} No story-specific notification is stated. Recommendations must respect resource privacy settings (\texttt{FR-27.5}, \texttt{NFR-PRIV-5}, \texttt{NFR-PRIV-7}).
\end{itemize}
\par\medskip

\hypertarget{sec:b11-source-alignment-notes}{}
\addcontentsline{toc}{subsubsection}{B11. Source Alignment Notes}
\subsubsection*{B11. Source Alignment Notes}
\par\medskip

\begin{itemize}
\item This use case is derived from US-26 in \texttt{userstory.md} section 8.2 and from \texttt{FR-26.1} to \texttt{FR-26.4} and \texttt{AC-26.1} to \texttt{AC-26.3} in the proposals.
\item Preference and relevance dependencies come from \texttt{FR-09.4}, \texttt{FR-09.5}, and \texttt{FR-25.5}; publication dependency comes from \texttt{FR-27.4}.
\item Recommendation segments, reason text, dismissal, refresh, and exact UI labels come from \texttt{library.md} as UI/UX recommendations; dismissal is analogous to \texttt{FR-24.4} but is not a confirmed US-26 FR.
\item \texttt{API.md} has no standalone library recommendation endpoint, so no library recommendation route, status code, or error code is invented.
\item Overlap with existing \texttt{UC-06-share-learning-resource.md}: UC-06 publishes resources into the library; UC-23 surfaces published resources as recommendations and saves them to a personal list.
\item No I-xx, E-xx, or G-xx model relationship ID is introduced; flow headings such as A3 and E1 are specification flow labels only.
\end{itemize}
\par\medskip

\par\medskip
\hypertarget{sec:ui-prototype-specimens}{}
\addcontentsline{toc}{subsection}{UI Prototype Specimens}
\subsection*{UI Prototype Specimens}
\par\medskip
\begin{figure}[H]
\centering
\includegraphics[width=0.75\textwidth,height=0.42\textheight,keepaspectratio]{"../../prototype/feed_algorithms_desktop.png"}
\par\vspace{2pt}
{\footnotesize \textbf{UI Prototype:} Feed: Algorithms - Desktop}
\end{figure}
\par\smallskip

\end{document}
